\documentclass{article}
\usepackage{tikz}
\usetikzlibrary{angles, quotes}
\usepackage{amssymb}
\usepackage{pgfplots}
\pgfplotsset{width=10cm,compat=1.9}
\usepgfplotslibrary{statistics}
\usepackage[left=2.54cm,right=2.54cm,top=2.54cm,bottom=2.54cm]{geometry}
\usepackage{tasks} %Paquete necesario para la producción de items horizontales para algunos ejercicios de tarea empleando el comando /begin{multicols}}{#}.../end{multicols} para ayudarnos a reducir las páginas
\usepackage{pdflscape} %comando que permite cambiar a orientación horizontal las hojas%
\usepackage{soul} %Paquete para permitir la compilación de acentos usando el comando {\hl{}}}\\
\sethlcolor{yellow(munsell)} %comando para definir el color para el subrayado usando el comando \hl
\usepackage[utf8]{inputenc}
\usepackage[spanish]{babel}
\usepackage{icomma}
\usepackage{ragged2e}
\usepackage{siunitx}
\usepackage{url}
\usepackage[colorlinks=true, urlcolor=blue,  linkcolor=black, citecolor=green]{hyperref}
\usepackage{pdfpages}
\usepackage{blindtext} %Paquete que genera texto ficticio (dummy text) con el comando: \blindtext
\usepackage{cancel} %in the preamble gives you four different modes of striking through: \cancel{text to cancel} draws a diagonal line (slash) through its argument, \bcancel{text to cancel} uses the negative slope (a backslash), \xcancel{text to cancel} draws an X (actually \cancel plus \bcancel), \cancelto{〈value〉}{〈expression〉} draws a diagonal arrow through the 〈expression〉pointing to the 〈value〉 (math-mode only)
\usepackage{csquotes}
\usepackage{afterpage}
\usepackage{parskip} 
\usepackage{float}
\usepackage{enumitem}
\usepackage{multicol}%Paquete que permite la creación aislada de columnas de texto. De esta forma se reduce la cantidad de páginas en nuestro documento
\usepackage{lipsum} %Paquete que genera texto de relleno al igual que \usepackage{blindtext}, se usa con el comando \lipsum[#-#] los #(hashtags) delimitan el número de páginas que deseamos ocupar del paquete lipsum para usarlos como relleno. 

% Definir el comando personalizado para la nota
\newcommand{\mynote}[1]{%
\begin{tikzpicture}[baseline=-0.75ex]
\node[draw, circle, fill=Apple Green!60, inner sep=2pt] (note) {\textbf{Nota:}};
\end{tikzpicture}%
\ \textit{#1}%
}

\newenvironment{Figura}
{\par\medskip\noindent\minipage{\linewidth}}
{\endminipage\par\medskip}
\usepackage{caption}
\usepackage[
backend=biber,
sorting=none,
url=true
]{biblatex} %bibliografía
\addbibresource{biblio.bib}
% Establecer el espacio entre las entradas de la bibliografía
\setlength{\bibitemsep}{1\baselineskip} % Puedes ajustar el valor según tus preferencias
\usepackage{amsmath,amsthm,amssymb,amsfonts}
\usepackage{pifont} %Permite la compilación del comando \xmark: es el símbolo de la tachita
\usepackage{mathtools} %permite la compilación de símbolos para matrices
\usepackage{empheq} %Paquete que se relaciona con \usepackage[most]{tcolorbox} para la creación de cuadros/cajas de colores para delimitar los resultados matemáticos o líneas de texto usando la línea de comando: \begin{empheq}[box={\mymath[colback=orange(webcolor)!70,drop lifted shadow]}]{equation*} \end{empheq}
\usepackage[most]{tcolorbox} %Paquete que se relaciona con \usepackage{empheq} para la creación de cuadros/cajas de colores para delimitar los resultados matemáticos o líneas de texto
\renewcommand{\qedsymbol}{$\blacksquare$}
\usetikzlibrary{positioning,decorations.pathreplacing} %Paquete que permite la compilación de llaves para cuadros sinópticos (brace diagram) 
\usepackage{schemata} %The schemata package is designed just for these brace diagrams
\usepackage{booktabs}

\newcommand\AB[2]{\schema{\schemabox{#1}}{\schemabox{#2}}} %comando o paquete necesario para crear cuadros sinópticos

\newtcbox{\mymath}[1][]{%
nobeforeafter, math upper, tcbox raise base,
enhanced, colframe=blue!30!black,
colback=blue!30, boxrule=1pt,
#1} %Comando importante para encasillar los resultados matemáticos en cajas de diferentes colores, se relaciona con los paquetes \usepackage{empheq} y \usepackage[most]{tcolorbox} 

\newenvironment{sysmatrix}[1]
{\left(\begin{array}{@{}#1@{}}}
{\end{array}\right)}
\newcommand{\ro}[1]{%
\xrightarrow{\mathmakebox[\rowidth]{#1}}%
}
\newlength{\rowidth}% row operation width
\AtBeginDocument{\setlength{\rowidth}{3em}}  %Comando importante para laproducción de lineas de operaciones entre matrices, método de Gauss o Gauss Jordan se relaciona con \newenvironment{sysmatrix}

%formato para cambiar el horario a español
\usepackage[spanish]{datetime2}
\DTMsetdatestyle{spanish}

\renewcommand{\today}{\DTMdisplaydate{\the\year}{\the\month}{\the\day }{-1}}
%%%%%%%%%%%%%%%%%%%%%%%%%%%%%%%%%%%%%%%%%

\usepackage{datetime}
\newcommand{\mycurrenttime}{\xxivtime}

%%%%%%%%%%%%%%%%%%%%%%%%%%%%%%%%%%%%%%%%%%%%%%%%%%%%%%%%%%%%%%%%%%%%%%%%%%%%%%%%%%%%%%%%%%%%%%%%%%%%%%%%%%%%%%%%%%%%%%%%%%%%%%%%%%%%%%%%%%%%%%%%%


%%%%%%%%%%%%%Caja de color para el título%%%%%%%%%%%%%%%%%%%%%%%%%%%%%%%%%%%%%%%%%%%%%%%%%%%%%%

\definecolor{myframecolor}{RGB}{1, 45, 75} % Prussian Blue
\definecolor{myboxcolor}{RGB}{0, 107, 92} % Apple Green

\newtcolorbox{mybox}{
enhanced,
colback=myboxcolor!7, % Color de fondo del cuadro
colframe=myframecolor, % Color del marco del cuadro
arc=0pt,
boxrule=1pt,
borderline west={2mm}{-10mm}{myframecolor}, % Borde en el lado izquierdo
sharp corners=southwest,
width=\linewidth,
before=\par\vspace{\bigskipamount}, % Espacio antes del cuadro
after=\par\vspace{\bigskipamount} % Espacio después del cuadro
}


\newcommand{\euler}{e} %Comando para producir letra e de euler
\newenvironment{remark}{\par\vfill\footnotesize % Vertical white space above the remark and smaller font size
\begin{list}{}{
	\leftmargin=80pt % Indentation on the left
	\rightmargin=60pt}\item\ignorespaces % Indentation on the right
\makebox[-2.5pt]{\begin{tikzpicture}[overlay]
		\node[draw=Horizon!60,line width=2.5pt,circle,fill=Horizon!25,font=\sffamily\bfseries,inner sep=4pt,outer sep=0pt] at (-19pt,5pt){\textcolor{Horizon}{Nota}};\end{tikzpicture}} % Blue Nota in a circle
\advance\baselineskip -1pt}{\end{list}\vskip5pt} % Tighter line spacing and white space after remark
\usepackage{graphicx}

\usepackage{titling}

\input{colores.tex}

%%Producir signos de cita textual``Asignatura''

\newcommand{\xmark}{\ding{55}} %%%comando que genera la tachita

\newenvironment{MyColorPar}[1]{%
\leavevmode\color{#1}\ignorespaces%
}{%
}%

%%%%%%%%%%%%%%%%%%% Título de la tarea, Nombre de alumno

\title{Clase \# 39 - Mecánica Teórica}
\author{\emph{JA} $\boldsymbol{\mid}$}
\date{\today}

%%%%%%%%%%%%%%%%%%% Datos de la Materia

\usepackage{fancyhdr}
\fancypagestyle{plain}{%  the preset of fancyhdr 
\fancyhf{} % clear all header and footer fields
\fancyfoot[R]{
\includegraphics[width=2cm]{LogoFCFMBUAP (1).png}}
\fancyfoot[L]{{\bfseries{\thedate{}}} a las {\bfseries{\mycurrenttime{}}} horas{} (GMT-6, H. Puebla de Zaragoza, Pue)}
\fancyhead[L]{Mecánica teórica}
\fancyhead[R]{\theauthor}
}
\makeatletter
\def\@maketitle{%
\newpage
\null
\vskip 1em%
\begin{center}%
\let \footnote \thanks
{\LARGE \@title \par}%
\vskip 1em%
%{\large \@date}%
\end{center}%
\par
\vskip 1em}
\makeatother

\usepackage{lipsum}  
\usepackage{cmbright}


%%%%%%%%%%%%%%%%%%%%%%%%%%%%%%%%%%%%%%%%%%%%%%%%%%%%%%%%%%%%%%%%%%%%%%%%%%%%%%%%%%%%%%%%%%%%%%%%%%%%%%%%%%%%%%%%%%%%%%%%%%%%%%%%%%%%%%%%%%%%%%%%%%%%%%%%%%%%%%%%%%%%%%%%%%%%%%%%%%%%%%%%%%%%%%%%%%%%%%

\begin{document}



\maketitle


%%%%%%%%%%%%%%%%%%% Datos del profesor, campus y aula

\begin{mybox}
\noindent\begin{tabular}{@{}ll}
	{\bfseries{Profesor}} & Dr. Gilberto Silva Ortigoza\\
	\textcolor{prussianblue}{\bfseries{Campus}} / \textcolor{trueblue}{\bfseries{Facultad}}  & \textcolor{prussianblue}{\bfseries{C.U. BUAP}} /  \textcolor{trueblue}{\bfseries{FCFM}} \\
	\textcolor{prussianblue}{\bfseries{Edificio}} /  \textcolor{trueblue}{\bfseries{EMA4}}    & 	\textcolor{prussianblue}{\bfseries{FM9}} /  \textcolor{trueblue}{\bfseries{302}}   
\end{tabular} 
\end{mybox} \vspace{1cm}

\tableofcontents

\vspace{2cm}

\section{Clase de hoy: Transformaciones canónicas, corchetes de Poisson y un poco de función generatriz}

\begin{figure}[H]
    \centering
    \includegraphics[width=0.9\textwidth]{1.png}
    \caption{ transformaciones canónicas.}
    \label{fig:canonica1}
\end{figure}

\begin{figure}[H]
    \centering
    \includegraphics[width=0.9\textwidth]{2.png}
    \caption{paréntesis de poisson de una transformación canónica.}
    \label{fig:canonica2}
\end{figure}

Cuando realizamos una transformación canónica, el \textbf{paréntesis de Poisson} entre las nuevas variables \((Q,P)\) debe conservar su valor igual a 1. Es decir:

\[
\left( 
\frac{\partial Q}{\partial q}\frac{\partial P}{\partial p}
- 
\frac{\partial Q}{\partial p}\frac{\partial P}{\partial q}
\right) = 1
\]

Este resultado expresa la \textbf{invariancia del paréntesis de Poisson} bajo transformaciones canónicas. En otras palabras, el Jacobiano de la transformación es igual a uno, lo que garantiza que el \textbf{elemento de volumen en el espacio fase permanece invariante}. 


\subsection{Invariancia}

Al igual que las \textbf{transformaciones de Lorentz} dejan invariante el elemento de línea en el espacio-tiempo, las transformaciones canónicas dejan invariante la estructura del espacio fase.  
Recordemos que las rotaciones son transformaciones que conservan la longitud de un vector. De manera análoga, en el espacio fase las coordenadas \((q,p)\) son solo etiquetas, mientras que lo físicamente significativo son los \textbf{observables}. 

Decir que un sistema es invariante a lo largo de un eje (por ejemplo, el eje $z$) implica que sus ecuaciones no cambian cuando nos desplazamos a lo largo de esa coordenada.

Así, las \textbf{transformaciones canónicas} son aquellas que se realizan en el espacio fase de modo que las ecuaciones de Hamilton \emph{mantienen su forma}.  
Del mismo modo, las \textbf{transformaciones de punto} son aquellas que dejan invariantes en forma las ecuaciones de Lagrange.


\subsection{Invariancia en forma de las ecuaciones}

Por ejemplo, en coordenadas cartesianas y polares las ecuaciones de Lagrange no tienen la misma apariencia algebraica, pero son equivalentes en contenido físico.  
Ser “invariante en forma” significa justamente eso: las ecuaciones conservan su estructura funcional. 

\[
\frac{d}{dt} \!\left( \frac{\partial L}{\partial \dot{q}_{i}} \right)
- \frac{\partial L}{\partial q_{i}} = 0
\qquad \Longleftrightarrow \qquad
\frac{d}{dt} \!\left( \frac{\partial L}{\partial \dot{s}_{i}} \right)
- \frac{\partial L}{\partial s_{i}} = 0
\]

De acuerdo con \textbf{Safko, Goldstein y Poole}, una transformación canónica es aquella que deja invariante la forma de las ecuaciones de Hamilton:

\[
\dot{q} = \frac{\partial H}{\partial p},
\qquad
\dot{p} = -\frac{\partial H}{\partial q}
\]

y bajo la transformación:

\[
\boxed{
Q = Q(q,p,t), \qquad P = P(q,p,t)
}
\]

no siempre se cumple que:

\[
\boxed{
\dot{Q} = \frac{\partial K}{\partial P}, 
\qquad
\dot{P} = -\frac{\partial K}{\partial Q}
}
\]

Por tanto, el criterio fundamental es que se conserve la estructura simpléctica, es decir, la \textbf{invariancia de los paréntesis de Poisson}.  
Según el Dr.~Ortigoza, una transformación canónica es aquella que deja invariante precisamente esta estructura de los paréntesis de poisson (\textcolor{Cinnabar}{\textbf{¡Checar las notas!}})


\subsection{Condición de canonicidad}

Si la transformación es canónica, entonces se cumplen los paréntesis de Poisson:

\[
[Q, P] = 1
\quad \Longleftrightarrow \quad
\frac{\partial Q}{\partial q}\frac{\partial P}{\partial p}
- 
\frac{\partial Q}{\partial p}\frac{\partial P}{\partial q} = 1
\]

Esta condición garantiza la existencia de una \textbf{función generatriz} asociada a la transformación.  
Es decir, si los paréntesis de Poisson se satisfacen, entonces existe una función generatriz $F$ que describe la transformación canónica.

\[
Q = Q(q,p,t),
\qquad
P = P(q,p,t)
\]


\subsection{Relación diferencial y función generatriz}

Podemos escribir la relación anterior de la siguiente manera:

\[
\frac{\partial}{\partial p}
\!\left(
P\frac{\partial Q}{\partial q} - p
\right)
=
\frac{\partial}{\partial q}
\!\left(
P\frac{\partial Q}{\partial p}
\right)
\]

De donde se obtiene:

\[
\frac{\partial P}{\partial p}\frac{\partial Q}{\partial q}
+ P\frac{\partial^{2} Q}{\partial p \partial q}
- 1
=
\frac{\partial P}{\partial q}\frac{\partial Q}{\partial p}
+ P\frac{\partial^{2} Q}{\partial q \partial p}
\]

Sea ahora una función \( w(q,p) \) tal que:

\[
dw = 
\boxed{
\frac{\partial w}{\partial q}\,dq
+ \frac{\partial w}{\partial p}\,dp
}
\]

La condición para que esta diferencial sea exacta es:

\[
\frac{\partial}{\partial p}
\!\left(
\frac{\partial w}{\partial q}
\right)
=
\frac{\partial}{\partial q}
\!\left(
\frac{\partial w}{\partial p}
\right)
\]

Por analogía, existe entonces una función \( F \) tal que:

\[
dF =
- \!\left( P\frac{\partial Q}{\partial q} - p \right)\!dq
- \!\left( P\frac{\partial Q}{\partial p} \right)\!dp
\]

Si esta expresión es exacta, significa que la derivada parcial del primer término con respecto a \( p \) es igual a la del segundo término con respecto a \( q \).  
Por tanto, si \( Q \) y \( P \) definen una transformación canónica, entonces existe una función \( F \) que cumple:

\[
\frac{\partial F}{\partial q} = -P\frac{\partial Q}{\partial q} + p,
\qquad
\frac{\partial F}{\partial p} = -P\frac{\partial Q}{\partial p}
\]

Estas ecuaciones indican que, si conocemos la transformación canónica, podemos obtener la función generatriz integrando dichas relaciones.  
Y de forma recíproca: si conocemos la función generatriz, podemos deducir completamente la transformación canónica.

\subsection{Como recordatorio el examen de kepler será el próximo martes 4 de noviebre}


\end{document}